\section{Canonical Gaussian domination with quantized transverse photons}
\label{pe-section}

This section extends the determinant argument to the full cutoff
Hamiltonian, including its quantized transverse photons. The photon
integration is justified first for a bounded coordinate truncation and
then removed by quadratic-form estimates. The reference energy is the
actual interacting neutral vacuum. No extensivity assumption about that
vacuum is used.

The sufficient restriction is
\begin{equation}
m_*:=m-\frac{13\alpha}{\pi^2}\Lambda_f>0.
\label{pe-mass-condition}
\end{equation}
It is a restriction on the bare parameters and the fixed fermionic
cutoff, not a mass-renormalization statement. The photon cutoff is any
finite \(\Lambda_\gamma\). The nuclear self-energy convention remains
that of Section~\ref{sec:model}. We now specify the photon regulator.
Take all transverse modes \(0<|k|\le\Lambda_\gamma\), two polarizations
\(\varepsilon_\lambda(k)\perp k\), and \(\omega_k=|k|\):
\begin{align}
H_\gamma&=\sum_{k,\lambda}\omega_k a^\dagger_{k,\lambda}a_{k,\lambda},
\\
A(y)&=\frac1{\sqrt V}\sum_{k,\lambda}
\frac{\varepsilon_\lambda(k)}{\sqrt{2\omega_k}}
\left(a_{k,\lambda}e^{ik\cdot y}
+a^\dagger_{k,\lambda}e^{-ik\cdot y}\right),\\
H_I&=e\int_{\T_L^3}:\!\psi^\dagger(y)\boldsymbol{\mathsf a}
\psi(y)\!:_P\cdot A(y)\,\dd y,\qquad e^2=4\pi\alpha .
\label{pe-field-definition}
\end{align}
Here \(y\) is spatial position; \(x\) below denotes photon coordinates.
The harmonic transverse \(k=0\) mode is explicitly fixed and omitted,
an additional convention not implied by Coulomb gauge alone.
There is no further infrared decoupling: the smallest retained
frequency, when present, is \(2\pi/L\). The Hilbert space is
\(\F\otimes\F_{\rm sym}(\mathfrak h_{\gamma,L})\).
We write \(H_0\) within this section for the full source-free Hamiltonian,
and \(E^{\vac}_{\Lambda,L}\) for its interacting neutral ground energy.
The current matrices on the finite fermion space are bounded, so
\(H_I\) is infinitesimally \(H_\gamma\)-bounded: for each retained mode,
\(\|a_k\Psi\|\le\omega_k^{-1/2}\|H_\gamma^{1/2}\Psi\|\) and
\(\|a_k^\dagger\Psi\|^2=\|a_k\Psi\|^2+\|\Psi\|^2\);
the scalar inequality \(\sqrt t\le\varepsilon t+(4\varepsilon)^{-1}\)
then gives the required relative bound after summing finitely many modes.
Kato--Rellich gives
self-adjointness on \(D(H_\gamma)\); the oscillator form estimates
below prove a lower bound and compact resolvent at each fixed regulator.

\subsection{Exact matching to the ordering algebra}
Set \(h_L=\hf\), \(C_L=C\), and
\(\kappa_L=\kappa_{\Lambda,L}\) from Section~\ref{sec:ordering}.
Lemmas~\ref{lem:wick} and \ref{lem:mass} give the exact full-field identity
\begin{equation}
H_0=H_\gamma+C_L+\mathrm d\Gamma(h_L)+H_I
+\frac\alpha{2V}\sum_{k\ne0}v_k\rho_{-k}\rho_k.
\label{pe-wick}
\end{equation}
In particular,
\begin{equation}
h_L=\sum_{j=1}^3\mathsf a_jv_j+\mathsf b M,
\qquad M\ge m_*I>0,\qquad \Tr h_L=0.
\label{pe-positive-mass}
\end{equation}
There are no additional ordering changes in this section.
For the scalar Hubbard--Stratonovich variables use
\(\widehat\rho(F)=\mathrm d\Gamma(F)-\Tr F/2\),
where \(F\) is Hermitian and spin scalar. This is the opposite
sign to physical charge density, absorbed into the source shifts.
The Coulomb operator is the finite sum of positive real squares
in \eqref{eq:squares}. Below \(d=4M_L\) and \(Q=N-d/2\).

\subsection{Positivity for arbitrary real photon-coordinate histories}

\begin{lemma}[Centered CAR trace with vector potentials and scalar kicks]
\label{pe-trace-positive}
Let
\[
h_j=\sum_{i=1}^3\mathsf a_i P_{i,j}+\mathsf b M_j,
\qquad M_j\ge m_*I>0,
\]
with Hermitian spatial coefficients. Let \(s_j\ge0\), with
\(\sum_js_j>0\), and let the kicks be arbitrary exponentials of
\(i\widehat\rho(F)\), with Hermitian spin-scalar \(F\), together
with an arbitrary charge twist \(e^{i\theta Q}\). The CAR trace of
the resulting ordered product of the free factors
\(e^{-s_jd\Gamma(h_j)}\) and these centered kicks is strictly positive.
\end{lemma}

\begin{proof}
Represent the one-particle product by an ODE
\(U'=-A(t)U\). On free intervals \(A=h_j\); on kick intervals
\(A=i\Phi(t)\), with \(\Phi\) Hermitian and spin scalar.
The charge twist is one more scalar kick. Equivalently the two types
can be combined as \(A=c(t)h(t)+i\Phi(t)\), with \(c\ge0\) and
positive total free time. Set
\(\mathsf C_5=\mathsf a_1\mathsf a_2\mathsf a_3\mathsf b\).
It is Hermitian and unitary, anticommutes with all the \(h_j\), and
commutes with every \(\Phi\). Consequently
\[
A^*\mathsf C_5+\mathsf C_5A=0,
\qquad U^*\mathsf C_5U=\mathsf C_5.
\]
The finite-dimensional CAR trace identity gives the centered trace as
\[
z=\exp\left(\frac12\int\operatorname{Tr}A(t)\,dt\right)
       \det(I+U).
\]
The exponential is precisely the scalar supplied by the centered kicks,
including \(e^{-i\theta d/2}\). Since \(\operatorname{Tr}h_j=0\),
\(\int\operatorname{Tr}A\) is purely imaginary. Pseudo-unitarity and
\(\det U=\exp(-\int\operatorname{Tr}A)\) therefore imply \(z=\bar z\).

This real trace cannot vanish. Otherwise \(Uv=-v\) for some nonzero
\(v\), and the corresponding ODE solution \(f\) is antiperiodic.
Taking the real part after multiplication by \(\mathsf b\) and
integration gives
\[
0=\operatorname{Re}\int f^*\mathsf b(\partial_t+A)f\,dt
  =\int c(t)f(t)^*M(t)f(t)\,dt>0.
\]
The derivative contributes only the vanishing endpoint difference.
The matrices \(\mathsf b\mathsf a_iP_{i,j}\) and
\(i\mathsf b\Phi\) are anti-Hermitian. The strict last inequality
uses invertibility of ODE propagation and positive total free time.
Finally deform all kicks to zero and deform the spatial matrices
\(P_{i,j}\) to zero and the positive \(M_j\) to \(m_*I\). The
argument excludes zeros throughout these continuous deformations.
At the endpoint the centered trace is
\(\det(I+e^{-\mathsf b m_*\sum_js_j})>0\), proving positivity.
\end{proof}

Real transverse vector potentials satisfy exactly the hypotheses of
this lemma: adding \(eP_f\boldsymbol{\mathsf a}\cdot A P_f\)
changes only the Hermitian \(P_{i,j}\), not \(M\). The proof needs
neither spatial locality of the momentum projection nor a photon gap
uniform in volume.

\subsection{Trace convergence for the bounded coordinate truncation}

Use standing-wave photon modes. Choose
\(\varepsilon_\lambda(-k)=\varepsilon_\lambda(k)\).
For each pair \(k,-k\), the unitary combinations
\(b_c=(a_k+a_{-k})/\sqrt2\) and
\(b_s=i(a_k-a_{-k})/\sqrt2\) make the vector potential a real
cosine/sine linear combination of \(b_c+b_c^\dagger\) and
\(b_s+b_s^\dagger\).
The equal frequencies preserve the oscillator Hamiltonian.
Its finite-mode Fock representation is therefore
\begin{equation}
H_\gamma=\frac12\sum_{a=1}^{d_\gamma}
  (-\partial_{x_a}^2+\omega_a^2x_a^2-\omega_a),
\qquad \omega_a>0,
\label{pe-oscillator}
\end{equation}
and the vector potential is linear in the real coordinates \(x_a\).
If no photon modes are retained the following photon integrals are
absent. Its Mehler heat kernel is strictly positive. The current
coupling is the multiplication operator
\(e\,d\Gamma(V(x))\), where
\(V(x)=P_f\boldsymbol{\mathsf a}\cdot A(x)P_f\).
There is no additional Wick scalar: \(\operatorname{Tr}(PV(x))=0\)
by the nonzero photon momenta and the symmetric free cutoff.

Choose a real continuous function \(0\le\chi_R\le1\), equal to one
on \(|x|\le R\) and zero on \(|x|\ge2R\). Replacing the coupling
by \(e\chi_R(x)d\Gamma(V(x))\) makes its matrix potential bounded.
Every truncated instantaneous one-body matrix still satisfies
\eqref{pe-positive-mass}.

\begin{lemma}[Trace-norm product formula used here]
\label{pe-trotter}
Let \(A\ge0\), with \(e^{-tA}\) trace class for every \(t>0\),
and let \(B_0,\ldots,B_J\) be bounded self-adjoint operators. Then,
for every \(b>0\),
\[
\left(e^{-bA/n}\prod_{j=0}^J e^{-bB_j/n}\right)^n
\longrightarrow e^{-b(A+\sum_jB_j)}
\quad\text{in trace norm}.
\]
The same holds after multiplication by a bounded charge projection.
\end{lemma}

\begin{proof}
For one bounded perturbation \(B\), expand both the product and the
bounded-perturbation Dyson series in powers of \(B\). If
\(s_0+\cdots+s_k=b\), Schatten H\"older gives
\[
\|e^{-s_0A}B e^{-s_1A}\cdots B e^{-s_kA}\|_1
\le\|B\|^k\operatorname{Tr}(e^{-bA}).
\]
The Schatten exponents are \(b/s_j\), with exponent infinity at
\(s_j=0\). The series tails are thus uniformly dominated by the
exponential series \(\operatorname{Tr}(e^{-bA})e^{b\|B\|}\).
At each fixed order, the terms with distinct product times are Riemann
sums for the time-ordered Dyson integral. The integrand is trace-norm
continuous in the interior of the time simplex; its boundary has
measure zero and the displayed bound controls boundary strips. Terms
with coincident times have total weight tending to zero. This proves
trace-norm convergence, since the Dyson series is the semigroup of the
bounded perturbation \(A+B\).

For several bounded \(B_j\), replacing
\(e^{-b\sum_jB_j/n}\) by \(\prod_je^{-bB_j/n}\) has operator-norm
error \(O(n^{-2})\), by the power series of bounded exponentials.
A telescoping sum, with the same Schatten H\"older estimate for the
intervening heat factors, bounds the total trace-norm error by
\(O(n^{-1})\operatorname{Tr}(e^{-bA})\) times a fixed exponential
in \(b\sum_j\|B_j\|\). Multiplication by a bounded projection
preserves convergence.
\end{proof}

\begin{proposition}[Canonical domination at bounded photon coordinates]
\label{pe-truncated-domination}
For the truncated coupling consider
\[
\mathcal H_R(\eta)=H_\gamma+C_L+
d\Gamma\bigl(h_L+e\chi_R(x)V(x)\bigr)
+\sum_{j=1}^J a_j\bigl(\widehat\rho(F_j)+\eta_j\bigr)^2,
\]
where \(a_j\ge0\), \(F_j\) are Hermitian spin-scalar matrices, and
\(\eta_j\in\mathbb R\). For every admissible integer \(q\) and
every inverse temperature \(b>0\),
\begin{equation}
\operatorname{Tr}_{Q=q}e^{-b\mathcal H_R(\eta)}
\le\operatorname{Tr}_{Q=0}e^{-b\mathcal H_R(0)}.
\label{pe-thermal-domination}
\end{equation}
\end{proposition}

\begin{proof}
Use Lemma~\ref{pe-trotter} with the oscillator as \(A\), the bounded
fermionic matrix potential as \(B_0\), and one bounded density square
for each remaining \(B_j\). At each product step use the exact scalar
Gaussian identity
\[
e^{-t a(X+\eta)^2}
=\frac1{\sqrt{2\pi}}\int_{\mathbb R}e^{-s^2/2}
          e^{i\sqrt{2ta}s(X+\eta)}\,ds.
\]
No commutation between different \(X\)'s is assumed. The separate
product factors handle that issue.

Write the canonical projection as
\(P_q=(2\pi)^{-1}\int_0^{2\pi}e^{i\theta(Q-q)}d\theta\).
The trace of each finite product is then an integral over oscillator
coordinates, scalar Gaussian variables, and \(\theta\). Its photon
weight is a cyclic product of positive Mehler kernels. After removing
the source phase and \(e^{-iq\theta}\), the remaining fermionic
integrand is strictly positive by Lemma~\ref{pe-trace-positive}.
Both removed factors have modulus one. The triangle inequality
therefore bounds the absolute value of the finite-product canonical
trace by the corresponding neutral, zero-source finite-product trace.

These integrals are legitimate ordinary finite-dimensional integrals.
The scalar kicks are unitary and the fermionic matrix potential is
bounded, so its trace is bounded uniformly in all Gaussian variables
by the finite CAR dimension times a fixed exponential. The cyclic
Mehler weight integrates to the finite oscillator partition function.
Lemma~\ref{pe-trotter} permits passage to the actual partition functions.
The left partition function is positive, giving
\eqref{pe-thermal-domination}.
\end{proof}

\subsection{Removal of the photon-coordinate truncation}

Let \(\mathcal H(\eta)\) denote the same Hamiltonian with
\(\chi_R=1\). At fixed volume and cutoffs there is a finite constant
\(C_1\) such that
\(\|e\,d\Gamma(V(x))\|\le C_1|x|\). Therefore
\begin{equation}
|\langle\Psi,(\mathcal H_R(\eta)-\mathcal H(\eta))\Psi\rangle|
\le\frac{C_1}{R}\langle\Psi,|x|^2\Psi\rangle.
\label{pe-form-error}
\end{equation}
Here is an explicit bound, including its charge-sector uniformity.
Write \(c_\gamma=\frac12\sum_a\omega_a\) and
\(\omega_*=\min_a\omega_a>0\). As quadratic forms,
\[
|x|^2\le\frac2{\omega_*^2}(H_\gamma+c_\gamma),\qquad
C_1|x|\le\frac{\omega_*^2}{4}|x|^2+\frac{C_1^2}{\omega_*^2}.
\]
Because all shifted density squares are positive, one may choose
\[
C_2=|C_L|+\|\mathrm d\Gamma(h_L)\|
+\frac{c_\gamma}{2}+\frac{C_1^2}{\omega_*^2}
\]
so that
\(\mathcal H_R(\eta),\mathcal H(\eta)\ge H_\gamma/2-C_2\).
This \(C_2\) is independent of \(R\), the charge sector, and the
source shifts \(\eta\) in the completed-square Hamiltonian.
Set \(\epsilon_R=2C_1/(R\omega_*^2)\) and
\(B=2C_2+c_\gamma\). Equation~\eqref{pe-form-error} gives
\begin{align}
\mathcal H_R(\eta)&\ge(1-2\epsilon_R)\mathcal H(\eta)-\epsilon_R B,\\
\mathcal H(\eta)&\ge(1-2\epsilon_R)\mathcal H_R(\eta)-\epsilon_R B.
\label{pe-explicit-squeeze}
\end{align}
For \(2\epsilon_R<1\), taking infima in each sector squeezes its
ground energy between quantities converging to the untruncated energy.
The common form domain is that of \(H_\gamma^{1/2}\), tensored with
the finite fermion space. Its compact embedding and the lower bound
give compact resolvent; the min--max principle bounds eigenvalues below
by those of \(H_\gamma/2-C_2\), giving trace-class heat semigroups.
If there are no photon modes, no truncation argument is needed.

First taking \(b\to\infty\) in \eqref{pe-thermal-domination}, and
then \(R\to\infty\), proves the exact ground-energy domination
\begin{equation}
\inf\sigma\bigl(\mathcal H(\eta)|_{Q=q}\bigr)
\ge\inf\sigma\bigl(\mathcal H(0)|_{Q=0}\bigr).
\label{pe-ground-domination}
\end{equation}
No limit uniform in \(L\) is required in this truncation argument:
the comparison has no truncation-dependent error before its removal,
and \eqref{pe-ground-domination} is then an exact statement at each
volume.

\subsection{One nucleus and the uniform many-nucleus bound}

Let \(D_L^{\rm band}\) be the Coulomb form restricted to
\(0<|k|\le2\Lambda_f\). Completing the density squares and using
\eqref{pe-ground-domination} yields, for any prescribed source \(\nu\),
\begin{equation}
\inf\sigma\bigl((H_0+\alpha D_L(\rho,\nu))|_{Q=q}\bigr)
\ge E^{\rm vac}_{\Lambda,L}
 -\frac\alpha2D_L^{\rm band}(\nu,\nu).
\label{pe-band-response}
\end{equation}
The reference on the right includes the photon interaction and all
Coulomb density squares. In particular, for a single point nucleus
of charge \(Z\) and the physical sector \(q=Z\),
\begin{equation}
E_{\Lambda,L}(R,Z)-E^{\rm vac}_{\Lambda,L}
\ge-\frac\alpha2\kappa_L Z^2
\ge-\frac{13\alpha}{\pi^2}\Lambda_f Z^2.
\label{pe-one-nucleus}
\end{equation}
This holds for every admissible sector and every \(L>0\).

For many nuclei, the direct band self-energy estimate alone should not
be used to discard the nuclear cross terms. Instead let \(p_t\) be
the periodic heat kernel and split
\begin{align*}
G_L^{\rm short}(x)&=4\pi\int_0^s p_t(x)\,dt-\frac{4\pi s}{V},\\
G_L^{\rm long}(x)&=\frac1V\sum_{k\ne0}
             \frac{4\pi e^{-s|k|^2}}{|k|^2}e^{ikx}.
\end{align*}
Both Fourier coefficient families are nonnegative. Split the
source-free density squares accordingly, and apply
\eqref{pe-ground-domination} with shifts only in the long-range
squares. The short-range source-free squares remain present throughout
the comparison. Thus the long-range source and nuclear terms together
are bounded below, relative to the same vacuum, by
\[
-\frac\alpha2G_L^{\rm long}(0)\sum_iZ_i^2.
\]
Indeed, completing the band squares costs at most
\(\alpha D_L^{\rm long}(\nu,\nu)/2\); the full long-range nuclear
cross terms cancel this quantity except for the individual diagonals.
Using the full long-range form only weakens the band estimate, because
the omitted Fourier tail is nonnegative.

For the short-range terms, \(\|\rho(x)\|\le2d_L\), where
\(d_L=M_L/V\), and
\(\|G_L^{\rm short}\|_{L^1}\le8\pi s\). Therefore their external
coupling is bounded below by \(-16\pi\alpha d_LsZ_{\rm tot}\).
Positivity of the heat kernel gives the nuclear lower bound
\[
\alpha\sum_{i<j}Z_iZ_jG_L^{\rm short}(R_i-R_j)
\ge-\frac{2\pi\alpha s}{V}Z_{\rm tot}^2
\ge-4\pi\alpha d_LsZ_{\rm tot},
\]
where the last step uses the exact neutrality capacity
\(Z_{\rm tot}\le2M_L\). Finally the same cubic-shell estimate as
above gives
\[
G_L^{\rm long}(0)
\le\frac{26}{\pi L}\sum_{j\ge1}
           e^{-4\pi^2sj^2/L^2}
\le\frac{13}{2\pi^{3/2}\sqrt s}.
\]
Combining the estimates proves
\begin{equation}
E_{\Lambda,L}(R,Z)-E^{\rm vac}_{\Lambda,L}
\ge-20\pi\alpha d_LsZ_{\rm tot}
 -\frac{13\alpha}{4\pi^{3/2}\sqrt s}\sum_iZ_i^2.
\label{pe-heat-bound}
\end{equation}

\begin{theorem}[Full cutoff relative stability in a specified cutoff window]
\label{pe-full-theorem}
Assume \eqref{pe-mass-condition}, \(L\ge\Lambda_f^{-1}\), finite
photon cutoff, distinct point nuclei, \(1\le Z_i\le Z_{\max}\),
and the exact neutral-sector admissibility condition
\(Z_{\rm tot}\le2M_L\). For the Hamiltonian of Section~\ref{sec:model}
with the photon terms \eqref{pe-field-definition}, including pair
creation and quantized photons,
\begin{equation}
E_{\Lambda,L}(R,Z)-E^{\rm vac}_{\Lambda,L}
\ge-\alpha\Lambda_f
\left[20\pi\left(1+\frac1\pi\right)^3Z_{\max}
 +\frac{13}{4\pi^{3/2}}Z_{\max}^2\right]K.
\label{pe-full-bound}
\end{equation}
There is no restriction on the photon number or on the number of pairs
other than the finite fermionic capacity.
\end{theorem}

\begin{proof}
Choose \(s=\Lambda_f^{-2}\) in \eqref{pe-heat-bound} and use
\[
d_L\le\left(\frac{\Lambda_f}{\pi}+\frac1L\right)^3
\le\left(1+\frac1\pi\right)^3\Lambda_f^3,
\qquad Z_{\rm tot}\le Z_{\max}K,
\qquad \sum_iZ_i^2\le Z_{\max}^2K.
\]
\end{proof}

Gaussian domination compares with the same interacting vacuum and
does not by itself prove its absolute extensivity.
It supplies neither a thermodynamic limit nor an ultraviolet limit.
Positive effective mass is the hypothesis of this proof; the displayed
linear cutoff restriction is only one sufficient criterion.
Later sections discuss other orderings and stronger criteria.
The result concerns the specified
fixed momentum cutoff and Wick prescription, and does not identify
physical renormalized masses or charges. Point-nucleus repulsion and
the fixed individual-self-energy subtraction have been preserved.
